% ==================== 章节4：具身智能 | 门把手、胶带和零号机 ====================

% --- 章节标题页（竖排） ---
\hirochapter{具身智能 | 门把手、胶带和零号机}

% --- 小剧场对话 ---
\section*{初星学园——2:00P.M. 活动准备室门口}
\phantomsection
\addcontentsline{toc}{section}{小剧场}

P：等一下，平板上的路线是对的。先把资料放到活动室，再去服装间拿水，最后回练习室。

広：嗯。很合理。

P：然后我们现在卡在门口。

広：因为门没有听从路线规划。

P：门把手不支持 API。

広：很遗憾。它比 agent 更接近现实。

P：你先别笑。资料快滑下去了。

広：我没有笑。只是觉得，今天的课题出现得很准确。

P：课题？

広：只会安排的东西，不能接手。它没有身体，所以门还是没有开。

P：所以你想给“小広计划”加一个身体。

広：不是一下子变成人。那样预算、伦理和普罗丢色桑的精神都会坏掉。

P：谢谢你还考虑了我的精神。

広：先做一个小的。会走，会停，会知道自己偏了。最好还能把水推过来一点。

P：低成本小车？

広：嗯。零号机。

P：听起来已经开始危险了。

広：不会太危险。很慢，很小，很笨。

P：这不是让人安心的形容。

広：但是很适合开始。

\clearpage

% --- 论文正文（IEEE风格，单栏） ---
\hiropapertitle{具身智能实践课：差速小车、SLAM 与两关节小臂}{広\textsuperscript{1}，Producer\textsuperscript{通讯}}{\textsuperscript{1}初星学园偶像学部\quad \textsuperscript{通讯}初星学园放学后代理计划}

\begin{abstract}
本章把具身智能从抽象概念落到一个小项目上：制作一台真实可做的低成本“小広试作机零号”。它不是人形机器人，也不是万能助手，而是一台差速轮小车，带有摄像头、轮速计、IMU 和一个简单的两关节小臂。本章围绕它在初星学园练习室里的三次失败，讲解机器人运动的最小数学、状态估计、SLAM、导航、局部避障、两关节机械臂和低层控制。重点不是把公式讲得华丽，而是看清一件事：机器人要接手现实任务，必须把“看见、估计、计划、执行、修正”连成闭环。
\end{abstract}
\keywords{具身智能、差速小车、状态估计、SLAM、导航、两关节机械臂}

\section{从“会说”到“会碰”}

前几章里，agent 已经学会了很多看起来像工作的事情：整理训练记录、安排路线、拆分任务、提醒 Producer 不要忘记带钥匙。
可是到了活动准备室门口，门拦住了agent的所有规划。

并非因为计划错了，计划甚至相当好。
问题在于，门不会因为计划合理就自动打开，水瓶不会因为表格写得漂亮就自己移动，快滑下来的资料也不会因为模型已经“理解任务”就停在半空。

P：那么今天的课题就是自动开门，移动水瓶和让资料漂浮？

広：嗯。今天的课题是具身智能呢，这些都可以做到的哦。

P：包括资料漂浮吗。。。

広：包括资料漂浮哦，如果加一些高速风扇的话。

具身智能关心的不是“智能体能不能给出答案”，而是答案能不能穿过身体，进入真实世界。
对机器人来说，任务不是停在一句“把水拿过来”上，而是要经过一串很麻烦的过程：

\begin{center}
观测世界 $\rightarrow$ 估计状态 $\rightarrow$ 规划路线 $\rightarrow$ 执行动作 $\rightarrow$ 根据反馈修正。
\end{center}

这个闭环听起来像机器人学，其实也很像偶像排练。
站位不是脑子里记住就够了，还要看镜子、听节拍、感觉脚下的地板、发现自己偏了，再把下一步拉回来。
只不过机器人没有“差不多”的身体直觉，它必须把每一件事拆成传感器、坐标、角度和控制量。

広：舞台上偏半步，会被发现。

P：小车偏半步，会撞椅子。

広：嗯。椅子比观众更严格。

\section{零号机任务书}

为了不让“小広计划”突然变成无法收场的大工程，Producer 给试作机画了一条很窄的边界。
它不是人形，不做拥抱，不扶人，不搬重物，也不在広体力不稳的时候靠近她。
它的第一阶段目标只有一个：

\begin{center}
从练习室门口出发，沿着胶带标出的区域移动到桌边，避开椅子，把水瓶推到安全胶带旁。
\end{center}

表 \ref{tab:zero-spec} 是零号机的最小配置。
这些部件并不神秘，甚至有点朴素。
但正因为朴素，它很适合当一台“能看见失败”的教学机器人。

\begin{table}[htbp]
\centering
\caption{小広试作机零号的最小配置}
\label{tab:zero-spec}
\small
\begin{tabular}{L{2.0cm} L{2.7cm} L{2.6cm}}
\toprule
模块 & 作用 & 本章关注的问题 \\
\midrule
差速轮底盘 & 前进、后退、原地转向 & 运动学、里程计误差 \\
摄像头 / 深度线索 & 看到胶带、桌腿、椅子、水瓶 & 视觉观测、遮挡、光照 \\
轮速计 / IMU & 估计移动距离和转向 & 状态估计、漂移修正 \\
两关节小臂 / 推杆 & 推近水瓶或轻道具 & 关节角、末端位置、力度 \\
急停与低速限制 & 防止撞人和撞坏东西 & 安全边界、回退策略 \\
\bottomrule
\end{tabular}
\end{table}

P：这个配置是不是太寒酸了？

広：很好。寒酸的话，失败会比较清楚。

P：你对“清楚”的要求总是有点残酷。

広：因为它现在最需要的不是威严。是知道自己在哪里。

\section{运动的最小数学}

\subsection{坐标系：先把练习室铺成一张纸}

要让零号机移动，第一步不是“让它聪明”，而是给它一个坐标系。
我们可以把练习室地板想象成一张平面图，取门口某个胶带角作为原点。
小车在这张图上的状态可以先写成
\begin{equation}
x_t =
\left[
\begin{array}{c}
p_x \\ p_y \\ \theta
\end{array}
\right]_t,
\end{equation}
其中 $p_x$ 和 $p_y$ 表示它在地板上的位置，$\theta$ 表示车头朝向。
这三个量听起来很少，但已经足够描述一台差速小车“在哪里、朝哪边”。

如果小车以线速度 $v_t$ 前进，以角速度 $\omega_t$ 转向，在很短时间 $\Delta t$ 内，它的位置可以近似更新为
\begin{equation}
\begin{array}{l}
p_{x,t+1} = p_{x,t} + v_t \cos \theta_t \Delta t,\\
p_{y,t+1} = p_{y,t} + v_t \sin \theta_t \Delta t,\\
\theta_{t+1} = \theta_t + \omega_t \Delta t.
\end{array}
\end{equation}

车头朝哪里，向前滚一小段以后，位置就会沿着那个方向变一点；同时，如果左右轮速度不一样，车头角度也会变。

P：所以它不是“走到桌边”，而是很多个很短的“往前一点、转一点”累积出来的。

広：对。偶像的走位也是这样。站到终点之前，先把每一步不要弄丢。

\subsection{差速轮：左右轮速度会决定转弯}

零号机用的是差速轮底盘。
左轮速度记作 $v_L$，右轮速度记作 $v_R$，两轮间距记作 $b$。
那么车体的近似线速度和角速度可以写成
\begin{equation}
v = \frac{v_R+v_L}{2}, \qquad
\omega = \frac{v_R-v_L}{b}.
\end{equation}

左右轮一样快，小车大致向前走。
右轮比左轮快，小车会向左转。
左轮比右轮快，小车会向右转。
如果一边前进一边后退，它甚至可以原地转身。

听起来很听话。
现实里没有这么乖。

轮子会轻微打滑，地板上贴胶带的位置摩擦不同，车轮半径也不一定完全一致。
于是公式里“应该走到”的地方，和真实世界里“实际停下”的地方，会慢慢分开。
这就是里程计漂移。

広：公式里的零号机很端正。

P：真实的零号机已经斜着进了椅子腿下面。

広：嗯。它开始像现实了。

\subsection{预测与观测：不要只相信自己走了多少}

机器人不能只靠轮子转了多少来判断自己在哪。
轮速计告诉它“我大概动了这么远”，摄像头告诉它“我看见胶带在画面左边”，IMU 告诉它“我刚才转了一个角度”。
这些信息都不是真相，只是线索。

一个常见的状态估计写法是
\begin{equation}
x_t = f(x_{t-1}, u_t) + w_t, \qquad
z_t = h(x_t, m) + v_t,
\end{equation}
其中 $x_t$ 是状态，$u_t$ 是控制输入，$z_t$ 是观测，$m$ 是地图或场地信息，$w_t$ 和 $v_t$ 分别代表运动噪声和观测噪声。

这两个式子可以翻成排练室语言：

\begin{itemize}[leftmargin=2em]
\item 先根据动作预测：我应该从这条胶带走到下一条胶带附近。
\item 再根据观测修正：可是摄像头看见桌腿比预期更近，所以我可能偏了。
\item 然后把偏差写回状态：下一步不要继续沿着错误方向走。
\end{itemize}

P：这和広训练时的复盘很像。

広：先预测能做到。然后失败。再修正。

P：你为什么说得像很开心。

広：因为这就是闭环。

\section{SLAM 与导航：练习室里的迷你地图}

\subsection{SLAM 要解决两件事}

SLAM 是 simultaneous localization and mapping，也就是同时定位与建图。
它要回答两个缠在一起的问题：

\begin{itemize}[leftmargin=2em]
\item 我在哪里？
\item 我周围的世界长什么样？
\end{itemize}

如果已经有一张精确地图，定位会容易很多。
如果已经知道自己在哪里，建图也会容易很多。
但零号机第一次进入练习室时，两件事都不完整。
它只能一边移动，一边把胶带、桌腿、椅子和墙角变成自己的迷你地图。

经典概率机器人学会把这个过程写成对后验分布的估计\cite{thrun2005prob}：
\begin{equation}
p(x_t, m \mid z_{1:t}, u_{1:t}).
\end{equation}
意思是：给定到目前为止的观测 $z_{1:t}$ 和控制输入 $u_{1:t}$，同时估计机器人轨迹 $x_t$ 与地图 $m$。

这类问题可以用滤波、图优化、视觉特征、惯性信息等多种方法解决。
VINS-Mono 用视觉和 IMU 做状态估计\cite{qin2018vinsmono}，ORB-SLAM3 则展示了视觉、视觉惯性和多地图 SLAM 的完整系统化路线\cite{campos2021orbslam3}。
在零号机这里，我们不需要真的实现论文级系统，但要理解它们共同面对的麻烦：每一步都有误差，误差会累积，累积到一定程度，小车就会把“桌边”认成“椅子旁边”。

\subsection{从地图到路线}

有了地图，还不等于能走。
导航至少要再拆成两层：

\begin{itemize}[leftmargin=2em]
\item 全局路径：从门口到桌边，大方向从哪里走。
\item 局部避障：眼前突然多了一把椅子，下一秒怎么绕。
\end{itemize}

在练习室里，全局路径像是排练前画在纸上的走位。
局部避障像是排练中有人临时把道具箱放错位置，你不能一边说“我的路线很正确”一边撞上去。

\begin{table}[htbp]
\centering
\caption{导航栈与初星学园里的对应物}
\label{tab:navi-map}
\small
\begin{tabular}{L{2.0cm} L{2.6cm} L{2.7cm}}
\toprule
机器人环节 & 学园里的对应物 & 问题 \\
\midrule
定位 & 站在舞台地标上 & 我在哪里 \\
建图 & 记住练习室障碍物 & 世界长什么样 \\
全局规划 & 走位路线 & 从哪里绕过去 \\
局部避障 & 临时让开椅子和同学 & 现实变了怎么办 \\
轨迹跟踪 & 听节拍后修正脚步 & 偏了怎么拉回来 \\
\bottomrule
\end{tabular}
\end{table}

局部避障方法有很多，Dynamic Window Approach 这类经典方法会在速度空间里挑选一段短时间内安全、可达、接近目标的速度\cite{fox1997dwa}。
直觉上，它不是一次性规划完整人生，而是在问：

\begin{center}
接下来这一小段，我能以什么速度走，既不撞东西，又离目标更近？
\end{center}

広：很像普罗丢色桑给我的训练菜单。

P：哪里像？

広：不是问“最终能不能完美”。是问“下一段能不能安全地难一点”。

P：这个解释我收下了。

\section{两关节小臂：碰到世界之前，先别用力过猛}

\subsection{为什么只做两关节}

零号机不是来表演杂技的。
如果第一阶段任务只是把水瓶推到安全胶带旁，一个两关节平面小臂就够了。
它可以理解成两个连杆接在一起：第一段长度 $l_1$，第二段长度 $l_2$；两个关节角分别是 $\theta_1$ 和 $\theta_2$。

末端位置可以写成
\begin{equation}
\begin{array}{l}
x = l_1 \cos \theta_1 + l_2 \cos(\theta_1+\theta_2),\\
y = l_1 \sin \theta_1 + l_2 \sin(\theta_1+\theta_2).
\end{array}
\end{equation}

这个式子是正运动学。
它回答的是：如果两个关节转成这样，末端会到哪里。
反过来，如果我们想让末端去某个点，就要解逆运动学。
本章不展开复杂求解，只记住两件事：

\begin{itemize}[leftmargin=2em]
\item 目标太远，超过 $l_1+l_2$，够不到。
\item 目标在可达范围内，也可能有“手肘朝上”和“手肘朝下”两种姿态。
\end{itemize}

P：所以它不是想推哪里就能推哪里。

広：对。身体的形状，会限制意志的范围。

P：这句话听起来突然很重。

広：但是很准确。

\subsection{小脑层：低层控制要快、稳、克制}

高层可以说“把水瓶推近一点”。
两关节小臂真正执行时，控制器要在很短时间里反复计算误差。
最朴素的形式可以写成 PID 控制：
\begin{equation}
u_t = K_p e_t + K_i \sum_{\tau=0}^{t} e_\tau + K_d(e_t-e_{t-1}),
\end{equation}
其中 $e_t$ 是目标和当前位置的误差，$u_t$ 是控制输出。
$K_p$ 像是“偏多少就拉多少”，$K_i$ 处理长期偏差，$K_d$ 则抑制动作太猛。

P：听起来像小脑。

広：嗯。它不负责想“为什么要拿水”。只负责别把水打翻。

如果执行层只追求快，就会冲过头。
如果执行层太软，就推不动。
接触任务里还需要一点顺应性：碰到水瓶时，不是硬撞，而是轻轻靠上去、感到阻力、再继续一点。
这也是为什么成熟机器人系统常把高层任务、路径规划和低层伺服分开。
高层慢一点没关系，低层必须及时。

\section{三次测试日志}

\subsection{第一次：能走，但找不到自己}

零号机第一次从门口出发时，路线看起来很顺利。
它沿着蓝色胶带走了约一米，车头轻轻偏向右侧。
Producer 以为这是小问题，直到它下一步把偏差继续放大，最后停在椅子腿前。

日志里，轮速计显示它“应该”还在胶带中央。
摄像头画面却显示胶带已经跑到画面左侧。
这就是只相信里程计的后果：每一步误差很小，但小误差会累积成明显偏航。

P：它不是突然走错的。

広：嗯。它是很认真地一点点错下去。

修正办法并不华丽：降低速度，提高观测频率，把摄像头检测到的胶带位置作为校正信号。
更完整的系统会把轮速计、IMU 和视觉观测放进滤波或图优化框架里。
零号机阶段先做一件小事就够了：不要只相信自己走了多少，要经常抬头看。

\subsection{第二次：能看，但不懂临时变化}

第二次测试前，Producer 故意把一把椅子挪到路线附近。
全局路径还是门口到桌边，地图上那条线没有问题。
问题出在局部。

零号机看见椅子晚了一点，停得也晚了一点。
它没有撞上去，但停在了一个尴尬的位置：再往前会碰椅子，往后退又会离开胶带。

P：它看见了，但反应太慢。

広：像我发现体力不够的时候。

P：你一般发现得太晚。

広：嗯。所以需要 Producer 急停。

这次失败说明，导航不是一张静态路线图。
练习室里的椅子会被挪，服装架会展开，路过的同学会临时经过。
因此局部规划要不断重算“接下来一点点能怎么走”。
在真实机器人里，这一层往往比高层语义更接近安全。
它不需要懂很多话，但必须知道“眼前这个东西不能撞”。

\subsection{第三次：能碰，但角度和力气不合适}

第三次，零号机终于接近桌边。
两关节小臂抬起来，末端对准水瓶。
从公式看，它已经把末端送到了目标附近。
从现实看，它的角度偏了几度。

于是水瓶不是被轻轻推近，而是被斜着碰了一下。
瓶身晃动，Producer 立刻按下急停。

広坐在旁边，看着水瓶停稳。
她没有先评价成功或失败，只问：

広：偏了多少？

P：末端大概偏了三厘米。力也大了。

広：嗯。很有价值。

P：你先确认水没洒。

広：也很有价值。

这次失败把机械臂的两个问题暴露得很清楚。
第一，末端位置不是只看目标点，还要看姿态和接触方向。
第二，接触不是单纯“到达”，而是要控制速度、力度和停止条件。
对零号机来说，下一步不是让它抓起水瓶，而是先学会在低速下轻轻推，并且一旦观测到水瓶晃动就停止。

\section{失败地图}

练习结束后，広在练习室地板上贴了三段不同颜色的胶带。
蓝色是偏航开始变明显的位置。
黄色是局部避障停得太晚的位置。
红色是水瓶晃动的位置。

P：你把失败点标得这么醒目，别人会以为这里发生过什么事故。

広：发生过。很小的事故。

P：你说得太平静了。

広：因为这是地图。

失败地图不是为了纪念失败，而是为了让下一次训练有入口。
零号机现在还不聪明。
它会偏，会犹豫，会把水瓶推得很难看。
但是它已经不再只是平板里的一段计划。
它开始把轮子、摄像头、角度和急停一起带进现实。

这也是具身智能最朴素的地方。
身体不是智能的外壳，而是题目本身的一部分。
只要进入现实，世界就会把所有“差不多”改成具体的误差：几厘米，几度，几毫秒，几次失败。

P：接下来呢？

広：让它先在一个小小的世界里多失败几次。

P：仿真？

広：嗯。再让它听懂一点人的话。

P：世界模型和 VLA？

広：下一章的目录。已经贴在地上了。

\begin{thebibliography}{99}
\bibitem{thrun2005prob} S. Thrun, W. Burgard, and D. Fox, \emph{Probabilistic Robotics}. Cambridge, MA, USA: MIT Press, 2005.
\bibitem{qin2018vinsmono} T. Qin, P. Li, and S. Shen, ``VINS-Mono: A Robust and Versatile Monocular Visual-Inertial State Estimator,'' \emph{IEEE Transactions on Robotics}, vol. 34, no. 4, pp. 1004--1020, 2018.
\bibitem{campos2021orbslam3} C. Campos, R. Elvira, J. J. Gomez Rodriguez, J. M. M. Montiel, and J. D. Tardos, ``ORB-SLAM3: An Accurate Open-Source Library for Visual, Visual-Inertial and Multi-Map SLAM,'' \emph{IEEE Transactions on Robotics}, vol. 37, no. 6, pp. 1874--1890, 2021.
\bibitem{fox1997dwa} D. Fox, W. Burgard, and S. Thrun, ``The Dynamic Window Approach to Collision Avoidance,'' \emph{IEEE Robotics \& Automation Magazine}, vol. 4, no. 1, pp. 23--33, 1997.
\end{thebibliography}
